% Token budgets and predicate-priority sampling: main-text section (#164).
% Replaces the draft's description of PPAS as a novel contribution and its
% future-work item "a genuine PPAS ablation" (done in #122/#128).
% D1 numbers: post-#137 population (726 known pairs), re-run at 66aabc3,
% artifacts/issue-164/reports/{ppas_ablation,ppas_controlled,ppas_validation}/.
% D5 (#122, be06539) and D2 (#128 follow-up, 4a42227) are unaffected by #137.
% Intervals: 95% paired bootstrap over the test reference sources of one pair
% (2,000 resamples, seed 42); no significance test.

\subsection{Token budgets and predicate-priority sampling}
\label{sec:ppas}

When an entity's key--value text (V3, V4, V6) exceeds the model's token
budget, the verbaliser keeps triples greedily in a fixed predicate-tier
order (for classes: lexical predicates such as labels, SKOS synonyms and
definitions, then type and hierarchy, then a few structural predicates;
predicates in no tier come last); we call this predicate-priority sampling
(PPAS). Only C3, C4, C5 and C7 can
invoke it, and at the budgets of the main conditions it rarely binds: in a
controlled re-implementation with M2's real tokenizer, no D1 text exceeds
420 tokens, and on D5 1.8\% of candidates do, with every selection policy
reaching the full-text MRR ($0.178$). The main results therefore say little
about PPAS, and we tested it separately.

\paragraph{Controlled ablation.}
One model (M2) and one key--value representation are held fixed and only the
selection policy varies, under outcome-blind stress budgets set before any
run (the pooled median full length, floored to a multiple of 8: 48 tokens on
D1, 80 on D5). On D1 at 48 tokens PPAS reaches MRR $0.494$, above random
selection ($0.406 \pm 0.001$ over five seeds), lexicographic deterministic
selection ($0.357$) and even the untruncated text ($0.453$). On D5 at 80
tokens PPAS cannot act: almost no D5 entity carries a tiered predicate, and
PPAS equals deterministic selection ($0.150$; random $0.136$, full $0.178$).

\paragraph{Eligibility, not priority.}
PPAS does two things at once: it orders the tiered predicates and, since
untiered predicates come last, it effectively excludes them under a tight
budget. Seven policies separate the two on D1 at 48 tokens. Restricting
selection to the tiered predicates explains the gain: with PPAS order in
both arms, so that the label is kept in each, it adds $0.192$
[$0.181$, $0.204$] MRR, and it adds $0.343$ [$0.330$, $0.355$] under
lexicographic order. With the eligible set fixed, PPAS order is slightly
\emph{worse} than lexicographic order ($-0.013$ [$-0.017$, $-0.010$]). The
pattern holds at 32 and 64 tokens (eligibility $+0.229$ and $+0.215$;
order $-0.014$ and $-0.002$) and on D2 at 48 tokens (eligibility $+0.093$
[$0.076$, $0.111$]; there the two orders select identical texts, so the
order effect is not identifiable). The apparent advantage of PPAS over
lexicographic order with every predicate eligible ($+0.137$
[$0.127$, $0.148$]) mostly reflects a weakness of the lexicographic
baseline: FMA and oboInOwl predicates sort before \texttt{rdfs:label}, so
under a tight budget lexicographic order drops the label of many
candidates, while an order that simply puts label, type and superclass
first, with no tiers, comes within $0.007$ of PPAS ($0.487$).

We therefore do not present PPAS's tier ordering as a contribution. Under a
tight budget, what matters is which predicates are eligible and that the
label survives. These results come from one model and one pair per dataset,
and we make no claim beyond them.
